% Part: methods
% Chapter: proofs
% Section: inference-patterns

\documentclass[../../../include/open-logic-section]{subfiles}

\begin{document}

\olfileid{mth}{prf}{pat}

\olsection{أنماط الاستدلال}

إنما يتركب البرهان من استدلالات مفردة، ولكل استدلال انتقال من معلوم
إلى لازم عنه. ويُنبَّه على هذا الانتقال بألفاظ كـ«إذن» و«وعليه»
و«لذلك». وقد يتكئ على حقيقة واحدة أو حقيقتين سبقتا في البرهان، سواء
كانتا من الفروض أو مما ثبت قبله. وربما رُمز إلى تلك الحقائق بحروف أو
أرقام، ثم أُحيل إليها عند استعمالها، ليكون مسار الحجة أبين. ويحسن
كذلك أن يُذكر \emph{وجه} لزوم النتيجة الجديدة عن مقدماتها. ولما كانت
جملة من صور الاستدلال كثيرة الدوران في البراهين، أفردنا الآتي لبيان
أشهرها؛ وأما الصور المركبة، كالاستقراء، فموضعها يأتي بعد.

وقد مر بنا منها ضرب واحد، وهو بسط التعريف أو إجراؤه. فبسط التعريف ردّ
الكلام المشتمل على المعرَّف إلى عبارة المُعرِّف. فإذا ثبت لنا، مثلًا،
في أثناء البرهان أن~$\OLId{latin-looped}{D} = \OLId{latin-looped}{E}$ (أ)، أجرينا تعريف~$=$ في المجموعات وقلنا:
«إذن، بالتعريف وبناء على (أ)، كل !!{element} من~$\OLId{latin-looped}{D}$ هو !!a{element}
من~$\OLId{latin-looped}{E}$، وصح العكس».

وقد يشتبه الأمر لأننا كثيرًا ما نورد مسوّغ الاستدلال قبل تمامه، لا
عند الفراغ منه. فافرض أن~$\OLId{latin-looped}{D} = \OLId{latin-looped}{E}$ لم تثبت بعد، وأن~$\OLId{latin-looped}{D} = \OLId{latin-looped}{E}$ عين
المطلوب. لنا أن نبسط الغاية بالتعريف فنقول: يلزم لإثبات~$\OLId{latin-looped}{D} = \OLId{latin-looped}{E}$ أن نثبت أن كل
!!{element} من~$\OLId{latin-looped}{D}$ هو !!a{element} من~$\OLId{latin-looped}{E}$، وأن العكس صحيح. ولو فُصّل
مسار البرهان لقيل: (أ) أثبت أن كل !!{element} من~$\OLId{latin-looped}{D}$ هو !!a{element}
من~$\OLId{latin-looped}{E}$؛ (ب) أثبت أن كل !!{element} من~$\OLId{latin-looped}{E}$ هو !!a{element} من~$\OLId{latin-looped}{D}$؛
(ج) يستنتج من (أ) و(ب)، بتعريف~$=$، أن~$\OLId{latin-looped}{D} = \OLId{latin-looped}{E}$. غير أن المعتاد جمع
ذلك في صورة أقرب إلى الآتية:
\begin{quote}
مرادنا أن نبيّن~$\OLId{latin-looped}{D} = \OLId{latin-looped}{E}$. وبحسب تعريف~$=$، يكافئ ذلك بيان أن كل
!!{element} من~$\OLId{latin-looped}{D}$ هو !!a{element} من~$\OLId{latin-looped}{E}$، وأن الأمر بالعكس كذلك.

(أ) \dots{} (برهان أن كل !!{element} من~$\OLId{latin-looped}{D}$ هو !!a{element} من~$\OLId{latin-looped}{E}$)
\dots

(ب) \dots{} (برهان أن كل !!{element} من~$\OLId{latin-looped}{E}$ هو !!a{element} من~$\OLId{latin-looped}{D}$)
\dots
\end{quote}

\subsection{استخدام اقتران}

وأقرب وجوه الاستدلال أن يؤخذ أحد جزأي الاقتران نتيجة. فإذا كان~$\OLId{latin-ordinary}{p}$
و$\OLId{latin-ordinary}{q}$ مفروضين أو ثابتين، جاز استخراج~$\OLId{latin-ordinary}{p}$، كما جاز استخراج~$\OLId{latin-ordinary}{q}$. وهذا
من البداهة بحيث يُطوى ذكره في أكثر المواضع. فبعد بسط تعريف~$\OLId{latin-looped}{D} = \OLId{latin-looped}{E}$،
مثلًا، نكون قد أثبتنا أن كل !!{element} من~$\OLId{latin-looped}{D}$ هو !!a{element} من
$\OLId{latin-looped}{E}$، وأن العكس صحيح؛ ومن مجموع القولين نأخذ أن كل !!{element} من~$\OLId{latin-looped}{E}$
هو !!a{element} من~$\OLId{latin-looped}{D}$، وهو المراد بقولنا «والعكس صحيح».

\subsection{إثبات اقتران}

وربما كان المطلوب اقترانًا، كأن يقال: «أثبت~$\OLId{latin-ordinary}{p}$ و$\OLId{latin-ordinary}{q}$». وحينئذ ينقسم
العمل قسمين: يثبت~$\OLId{latin-ordinary}{p}$ أولًا، ثم يثبت~$\OLId{latin-ordinary}{q}$. ولزيادة الإيضاح يمكن أن
تجعل في مسودة البرهان عنوانين: «(\OLClassicalDigits{1}) أثبت~$\OLId{latin-ordinary}{p}$» في صدر الصفحة، و«(\OLClassicalDigits{2})
أثبت~$\OLId{latin-ordinary}{q}$» في وسطها. وقد لا يُصرح في السؤال بأنه اقتران، ثم يكشفه فك
تعريف. فإذا كان المطلوب، مثلًا،~$\OLId{latin-looped}{D} = \OLId{latin-looped}{E}$، آل تعريف~$=$ إلى إثبات أن
كل !!{element} من~$\OLId{latin-looped}{D}$ هو !!a{element} من~$\OLId{latin-looped}{E}$ \emph{و}أن كل
!!{element} من~$\OLId{latin-looped}{E}$ هو !!a{element} من~$\OLId{latin-looped}{D}$.

\subsection{إثبات فصل}

وأما إذا كان المطلوب فصلًا، أي قولًا من نحو «$\OLId{latin-ordinary}{p}$ أو~$\OLId{latin-ordinary}{q}$»، فكفاية
البرهان في إظهار صدق أحد الطرفين. غير أن الفروض قلما تعيّن أحدهما
مباشرة؛ والأغلب أن تكون الفروض نفسها فصلية، أو أن يراد إثبات أن كل
شيء من صنف ما يتصف بإحدى صفتين، مع اختلاف الأفراد في أيتهما تتحقق.
وعندئذ يكون البرهان بالحالات، الآتي بيانه، هو الطريق الملائم.

\subsection{البرهان الشرطي}

وكثير من المبرهنات شرطيات، أي المطلوب فيها بيان أن تحقق~$\OLId{latin-ordinary}{p}$ يستلزم
صدق~$\OLId{latin-ordinary}{q}$. وطريق إعدادها بيّن: تُفرض مقدمة الشرطية، وهي~$\OLId{latin-ordinary}{p}$ هنا، ثم
يُستخرج منها التالي~$\OLId{latin-ordinary}{q}$. فإذا كان نص المبرهنة «إذا~$\OLId{latin-ordinary}{p}$ فإن~$\OLId{latin-ordinary}{q}$»، كان
أول البرهان «افترض~$\OLId{latin-ordinary}{p}$»، وكان آخره بلوغ~$\OLId{latin-ordinary}{q}$.

وتختلف العبارة عن الشرطية ولا يختلف معناها؛ فقد يقال بدل «إذا~$\OLId{latin-ordinary}{p}$
فإن~$\OLId{latin-ordinary}{q}$»: «$\OLId{latin-ordinary}{p}$ فقط إذا~$\OLId{latin-ordinary}{q}$»، أو «$\OLId{latin-ordinary}{q}$ إذا~$\OLId{latin-ordinary}{p}$»، أو «$\OLId{latin-ordinary}{q}$، بشرط~$\OLId{latin-ordinary}{p}$».
وفي الجميع يُفرض~$\OLId{latin-ordinary}{p}$ ويُثبت منه~$\OLId{latin-ordinary}{q}$. وأما التكافؤ، أعني «$\OLId{latin-ordinary}{p}$ إذا
وفقط إذا~$\OLId{latin-ordinary}{q}$»، فمركب من شرطيتين: إذا~$\OLId{latin-ordinary}{p}$ فإن~$\OLId{latin-ordinary}{q}$، وإذا~$\OLId{latin-ordinary}{q}$ فإن~$\OLId{latin-ordinary}{p}$؛
فيعاد البرهان الشرطي مرة لكل جهة. وقد يُثبت التكافؤ أحيانًا بسلسلة
من عبارات «إذا وفقط إذا» أولها «$\OLId{latin-ordinary}{p}$» وآخرها «$\OLId{latin-ordinary}{q}$»، بشرط التحقق من أن
كل حلقة في السلسلة تكافؤ حقًّا، لا استلزام من جهة واحدة.

\subsection{الادعاءات الكلية}

أما استعمال الدعوى الكلية فقاعدته أن ما ثبت لكل فرد ثبت للفرد
المعيّن. فإذا كان من فروض البرهان~$\OLId{latin-looped}{A} \subseteq \OLId{latin-looped}{B}$، فمعناه، بعد بسط
$\subseteq$، أن كل~$\OLId{latin-ordinary}{x} \in \OLId{latin-looped}{A}$ يحقق~$\OLId{latin-ordinary}{x} \in \OLId{latin-looped}{B}$. فإن عُلم بعد ذلك أن
$\OLId{latin-ordinary}{z} \in \OLId{latin-looped}{A}$، لزم~$\OLId{latin-ordinary}{z} \in \OLId{latin-looped}{B}$.

وقد يبدو إثبات ادعاء كلي عسيرًا قليلًا. وتأخذ هذه العبارات عادة
الصورة الآتية: «إذا كانت لـ$\OLId{latin-ordinary}{x}$ الخاصية~$\OLId{latin-looped}{P}$، فله الخاصية~$\OLId{latin-looped}{Q}$»، أو «كل
$\OLId{latin-looped}{P}$ هو~$\OLId{latin-looped}{Q}$». وقد لا تطابق هذه الصورة تمامًا بالطبع، وتحتاج إلى شيء
من المران لتعرف المطلوب إثباته على وجه الدقة. لكننا كثيرًا ما نضطر
إلى إثبات أن جميع الأشياء ذات خاصية ما لها خاصية معينة أخرى.

وطريقة إثبات ادعاء كلي هي إدخال أسماء أو متغيرات للأشياء التي لها
الخاصية الأولى، ثم بيان أن لها الخاصية الأخرى أيضًا. وقد نصوغ ذلك
بالقول إنه لكي تثبت شيئًا لـ\emph{كل}~$\OLId{latin-looped}{P}$، عليك إثباته لـ$\OLId{latin-looped}{P}$
\emph{اعتباطي}. والاسم المدخل اسم لـ$\OLId{latin-looped}{P}$ اعتباطي. ونستخدم عادة حروفًا
مفردة أسماء لهذه الأشياء الاعتباطية، وتتبع الحروف في العادة أعرافًا:
فنستخدم مثلًا~$\OLId{latin-ordinary}{n}$ للأعداد الطبيعية، و$!A$ لـ!!{formula}s، و$\OLId{latin-looped}{A}$
للمجموعات، و$\OLId{latin-ordinary}{f}$ للدوال، وهلم جرًا.

والحيلة هي الحفاظ على العمومية طوال البرهان. تبدأ بافتراض أن شيئًا
اعتباطيًا («$\OLId{latin-ordinary}{x}$») له الخاصية~$\OLId{latin-looped}{P}$، وتبيّن (استنادًا فقط إلى التعريفات
أو إلى ما يجوز لك افتراضه) أن لـ$\OLId{latin-ordinary}{x}$ الخاصية~$\OLId{latin-looped}{Q}$. ولأنك لم تحدد ما
$\OLId{latin-ordinary}{x}$ على وجه الخصوص سوى أن له الخاصية~$\OLId{latin-looped}{P}$، يمكنك أن تقرر أن كل شيء
له~$\OLId{latin-looped}{P}$ يملك الخاصية~$\OLId{latin-looped}{Q}$. وباختصار، ينوب~$\OLId{latin-ordinary}{x}$ عن \emph{جميع} الأشياء
ذات الخاصية~$\OLId{latin-looped}{P}$.

\begin{prop}
  لكل مجموعتين~$\OLId{latin-looped}{A}$ و$\OLId{latin-looped}{B}$، يكون~$\OLId{latin-looped}{A} \subseteq \OLId{latin-looped}{A} \cup \OLId{latin-looped}{B}$.
\end{prop}

\begin{proof}
  خذ~$\OLId{latin-looped}{A}$ و$\OLId{latin-looped}{B}$ مجموعتين كيف اتفق. والمطلوب
  $\OLId{latin-looped}{A} \subseteq \OLId{latin-looped}{A} \cup \OLId{latin-looped}{B}$؛ وهذا، بتعريف~$\subseteq$، أن يكون لكل
  $\OLId{latin-ordinary}{x}$: إذا كان~$\OLId{latin-ordinary}{x} \in \OLId{latin-looped}{A}$ كان~$\OLId{latin-ordinary}{x} \in \OLId{latin-looped}{A} \cup \OLId{latin-looped}{B}$. فليكن~$\OLId{latin-ordinary}{x}$
  عنصرًا اعتباطيًّا من~$\OLId{latin-looped}{A}$، ولنفرض~$\OLId{latin-ordinary}{x} \in \OLId{latin-looped}{A}$؛ ونريد إثبات
  $\OLId{latin-ordinary}{x} \in \OLId{latin-looped}{A} \cup \OLId{latin-looped}{B}$. ولما كان~$\OLId{latin-ordinary}{x} \in \OLId{latin-looped}{A}$، صدق أن~$\OLId{latin-ordinary}{x} \in \OLId{latin-looped}{A}$ أو
  $\OLId{latin-ordinary}{x} \in \OLId{latin-looped}{B}$. ومن ثم~$\OLId{latin-ordinary}{x} \in \Setabs{\OLId{latin-ordinary}{x}}{\OLId{latin-ordinary}{x} \in \OLId{latin-looped}{A} \lor \OLId{latin-ordinary}{x} \in \OLId{latin-looped}{B}}$،
  وهو، بحسب تعريف~$\cup$، عين القول إن~$\OLId{latin-ordinary}{x} \in \OLId{latin-looped}{A} \cup \OLId{latin-looped}{B}$.
\end{proof}


\subsection{البرهان بالحالات}

افترض أن لديك فصلًا بوصفه افتراضًا أو نتيجة مثبتة بالفعل---أي إنك
افترضت أو أثبتت صدق~$\OLId{latin-ordinary}{p}$ أو~$\OLId{latin-ordinary}{q}$. وتريد إثبات~$\OLId{latin-ordinary}{r}$. وتفعل ذلك في
خطوتين: تفترض أولًا صدق~$\OLId{latin-ordinary}{p}$ وتثبت~$\OLId{latin-ordinary}{r}$، ثم تفترض صدق~$\OLId{latin-ordinary}{q}$ وتثبت~$\OLId{latin-ordinary}{r}$
مرة أخرى. وينجح ذلك لأننا نفترض أو نعلم أن أحد البديلين متحقق.
وتثبت الخطوتان أن أيًا منهما يكفي لصدق~$\OLId{latin-ordinary}{r}$. (وإذا صدقا معًا، لم يكن
لدينا سبب واحد، بل سببان، لصدق~$\OLId{latin-ordinary}{r}$. ولا يلزم أن نثبت على حدة صدق~$\OLId{latin-ordinary}{r}$
بافتراض~$\OLId{latin-ordinary}{p}$ و$\OLId{latin-ordinary}{q}$ معًا.) ولكي نشير إلى ما نفعله، نعلن أننا «نميز
الحالات». افترض مثلًا أننا نعلم~$\OLId{latin-ordinary}{x} \in \OLId{latin-looped}{B} \cup \OLId{latin-looped}{C}$. وتُعرّف
$\OLId{latin-looped}{B} \cup \OLId{latin-looped}{C}$ بأنها~$\Setabs{\OLId{latin-ordinary}{x}}{\OLId{latin-ordinary}{x} \in \OLId{latin-looped}{B} \text{ أو } \OLId{latin-ordinary}{x} \in \OLId{latin-looped}{C}}$. أي
إن~$\OLId{latin-ordinary}{x} \in \OLId{latin-looped}{B}$ أو~$\OLId{latin-ordinary}{x} \in \OLId{latin-looped}{C}$ بحسب التعريف. ونثبت من ذلك~$\OLId{latin-ordinary}{x} \in \OLId{latin-looped}{A}$
بافتراض~$\OLId{latin-ordinary}{x} \in \OLId{latin-looped}{B}$ أولًا وإثبات~$\OLId{latin-ordinary}{x} \in \OLId{latin-looped}{A}$ من هذا الافتراض، ثم
بافتراض~$\OLId{latin-ordinary}{x} \in \OLId{latin-looped}{C}$ وإثبات~$\OLId{latin-ordinary}{x} \in \OLId{latin-looped}{A}$ مرة أخرى من هذا الافتراض. وتكتب
«نميز الحالات» تحت الافتراض، ثم «الحالة (\OLClassicalDigits{1}):~$\OLId{latin-ordinary}{x} \in \OLId{latin-looped}{B}$» تحتها،
و«الحالة (\OLClassicalDigits{2}):~$\OLId{latin-ordinary}{x} \in \OLId{latin-looped}{C}$» في منتصف الصفحة. ثم تملأ النصف الأعلى
والنصف الأسفل من الصفحة.

ويشتد نفع هذا الطريق إذا كان المطلوب نفسه فصلًا. ومن أبسط أمثلته:

\begin{prop}
افترض~$\OLId{latin-looped}{B} \subseteq \OLId{latin-looped}{D}$ و$\OLId{latin-looped}{C} \subseteq \OLId{latin-looped}{E}$. عندئذ
$\OLId{latin-looped}{B} \cup \OLId{latin-looped}{C} \subseteq \OLId{latin-looped}{D} \cup \OLId{latin-looped}{E}$.
\end{prop}

\begin{proof}
  افترض (أ)~$\OLId{latin-looped}{B} \subseteq \OLId{latin-looped}{D}$ و(ب)~$\OLId{latin-looped}{C} \subseteq \OLId{latin-looped}{E}$. بحسب التعريف، كل
  $\OLId{latin-ordinary}{x} \in \OLId{latin-looped}{B}$ يحقق أيضًا~$\OLId{latin-ordinary}{x} \in \OLId{latin-looped}{D}$ (ج)، وكل~$\OLId{latin-ordinary}{x} \in \OLId{latin-looped}{C}$ يحقق أيضًا
  $\OLId{latin-ordinary}{x} \in \OLId{latin-looped}{E}$ (د). ولكي نبيّن~$\OLId{latin-looped}{B} \cup \OLId{latin-looped}{C} \subseteq \OLId{latin-looped}{D} \cup \OLId{latin-looped}{E}$، علينا أن
  نبيّن أنه إذا كان~$\OLId{latin-ordinary}{x} \in \OLId{latin-looped}{B} \cup \OLId{latin-looped}{C}$، فإن~$\OLId{latin-ordinary}{x} \in \OLId{latin-looped}{D} \cup \OLId{latin-looped}{E}$ (بحسب
  تعريف~$\subseteq$). ويكون~$\OLId{latin-ordinary}{x} \in \OLId{latin-looped}{B} \cup \OLId{latin-looped}{C}$ إذا وفقط إذا كان
  $\OLId{latin-ordinary}{x} \in \OLId{latin-looped}{B}$ أو~$\OLId{latin-ordinary}{x} \in \OLId{latin-looped}{C}$ (بحسب تعريف~$\cup$). وبالمثل، يكون
  $\OLId{latin-ordinary}{x} \in \OLId{latin-looped}{D} \cup \OLId{latin-looped}{E}$ إذا وفقط إذا كان~$\OLId{latin-ordinary}{x} \in \OLId{latin-looped}{D}$ أو~$\OLId{latin-ordinary}{x} \in \OLId{latin-looped}{E}$. ومن ثم
  علينا أن نبيّن: لكل~$\OLId{latin-ordinary}{x}$، إذا كان~$\OLId{latin-ordinary}{x} \in \OLId{latin-looped}{B}$ أو~$\OLId{latin-ordinary}{x} \in \OLId{latin-looped}{C}$، فإن
  $\OLId{latin-ordinary}{x} \in \OLId{latin-looped}{D}$ أو~$\OLId{latin-ordinary}{x} \in \OLId{latin-looped}{E}$.

  \begin{quote}
  وما صنعناه إلى الآن إلا بسط التعريفات!{} فقد عبّرنا عن القضية من
  غير~$\subseteq$ و$\cup$، وبقي إثبات دعوى كلية شرطية. وعلى ما تقدم،
  نأخذ~$\OLId{latin-ordinary}{x}$ اعتباطيًّا ونفرض فيه صدق طرف «إذا»، ثم نبيّن له طرف «فإن».
  أي نفرض~$\OLId{latin-ordinary}{x} \in \OLId{latin-looped}{B}$ أو~$\OLId{latin-ordinary}{x} \in \OLId{latin-looped}{C}$ ونثبت~$\OLId{latin-ordinary}{x} \in \OLId{latin-looped}{D}$ أو
  $\OLId{latin-ordinary}{x} \in \OLId{latin-looped}{E}$.\footnote{هذه الفقرة شرح لطريقة العمل، وليست جزءًا من
  البرهان نفسه؛ فلا يلزم استيفاء هذا التفصيل كلما كتبت برهانًا.}
  \end{quote}

  افرض~$\OLId{latin-ordinary}{x} \in \OLId{latin-looped}{B}$ أو~$\OLId{latin-ordinary}{x} \in \OLId{latin-looped}{C}$. والمطلوب~$\OLId{latin-ordinary}{x} \in \OLId{latin-looped}{D}$ أو
  $\OLId{latin-ordinary}{x} \in \OLId{latin-looped}{E}$؛ فنقسم إلى حالتين.

  الحالة \OLClassicalDigits{1}:~$\OLId{latin-ordinary}{x} \in \OLId{latin-looped}{B}$. فمن (ج)،~$\OLId{latin-ordinary}{x} \in \OLId{latin-looped}{D}$؛ ومن ثم~$\OLId{latin-ordinary}{x} \in \OLId{latin-looped}{D}$ أو
  $\OLId{latin-ordinary}{x} \in \OLId{latin-looped}{E}$. (وهذا هو الاستدلال الذي مر في القسم الفرعي السابق.)

  الحالة \OLClassicalDigits{2}:~$\OLId{latin-ordinary}{x} \in \OLId{latin-looped}{C}$. فمن (د)،~$\OLId{latin-ordinary}{x} \in \OLId{latin-looped}{E}$؛ ومن ثم~$\OLId{latin-ordinary}{x} \in \OLId{latin-looped}{D}$ أو
  $\OLId{latin-ordinary}{x} \in \OLId{latin-looped}{E}$.
 \end{proof}


\subsection{إثبات ادعاء وجودي}

إذا كان المطلوب دعوى وجودية، جاء السؤال غالبًا على نحو «أثبت وجود
$\OLId{latin-ordinary}{x}$ بحيث~$\dots \OLId{latin-ordinary}{x} \dots$»؛ أي أثبت أن في الأشياء ما يتصف بالخاصية
«$\dots \OLId{latin-ordinary}{x} \dots$». فلا بد حينئذ من تعيين شاهد صالح ثم إثبات الخاصية
له. وإن بدا الأمر قريبًا، فقد يعسر هذا النوع من البراهين؛ إذ يقتضي
في العادة \emph{بناء} شيء أو \emph{تعريفه}، ثم إقامة الدليل على أن
المبني أو المعرّف يفي بالمطلوب. وربما عسر العثور على الشاهد، أو إثبات
صفته، بل ربما دقّ البرهان على أن التعريف عيّن شيئًا أصلًا!{}

ويبدأ ذلك عادة بقولك: «ليكن~$\OLId{latin-ordinary}{x}$ هو \dots»، فتجعل~$\dots$ وصفًا لما
تقصد. وقد يلزمك أولًا إثبات أن هذا الوصف يتحقق في شيء موجود، ثم تمضي
إلى بيان أن لـ$\OLId{latin-ordinary}{x}$ الخاصية~$\OLId{latin-looped}{Q}$. وهذا مثال يسير.

\begin{prop}
  افترض~$\OLId{latin-ordinary}{x} \in \OLId{latin-looped}{B}$. عندئذ توجد~$\OLId{latin-looped}{A}$ بحيث~$\OLId{latin-looped}{A} \subseteq \OLId{latin-looped}{B}$ و
  $\OLId{latin-looped}{A} \neq \emptyset$.
\end{prop}

\begin{proof}
  افرض~$\OLId{latin-ordinary}{x} \in \OLId{latin-looped}{B}$، واجعل~$\OLId{latin-looped}{A} = \{\OLId{latin-ordinary}{x}\}$.
  \begin{quote}
    عرّفنا هنا المجموعة~$\OLId{latin-looped}{A}$ بتعداد !!{element}sها. ولأننا نفترض أن
    $\OLId{latin-ordinary}{x}$ شيء، ولأن بوسعنا دائمًا تكوين مجموعة بتعداد !!{element}sها،
    فلا يلزم أن نبيّن هنا أننا نجحنا في تعريف مجموعة~$\OLId{latin-looped}{A}$. لكن ما زال
    علينا أن نبيّن أن لـ$\OLId{latin-looped}{A}$ الخواص التي تتطلبها القضية. ولا يكتمل
    البرهان من دون ذلك!{}
  \end{quote}
  ولما كان~$\OLId{latin-ordinary}{x} \in \OLId{latin-looped}{A}$، لزم~$\OLId{latin-looped}{A} \neq \emptyset$.
  \begin{quote}
    ووجه ذلك تعريف~$\OLId{latin-looped}{A}$ بأنها~$\{\OLId{latin-ordinary}{x}\}$، مع الحقيقتين الظاهرتين
    $\OLId{latin-ordinary}{x} \in \{\OLId{latin-ordinary}{x}\}$ و$\OLId{latin-ordinary}{x} \notin \emptyset$.
  \end{quote}
  ثم إن~$\OLId{latin-ordinary}{x}$ هو !!{element} الوحيد من~$\{\OLId{latin-ordinary}{x}\}$، ومع~$\OLId{latin-ordinary}{x} \in \OLId{latin-looped}{B}$ يكون كل
  !!{element} من~$\OLId{latin-looped}{A}$ !!a{element} من~$\OLId{latin-looped}{B}$ أيضًا. فينتج، بتعريف
  $\subseteq$، أن~$\OLId{latin-looped}{A} \subseteq \OLId{latin-looped}{B}$.
\end{proof}

\subsection{استخدام الادعاءات الوجودية}

افترض أنك تعرف صدق ادعاء وجودي ما (لأنك أثبته، أو لأنه فرضية يجوز لك
استخدامها)، وليكن «لبعض~$\OLId{latin-ordinary}{x}$،~$\OLId{latin-ordinary}{x} \in \OLId{latin-looped}{A}$»، أو «يوجد~$\OLId{latin-ordinary}{x} \in \OLId{latin-looped}{A}$». فإذا
أردت استخدامه في برهانك، أمكنك ببساطة أن تتظاهر بأن لديك اسمًا لأحد
الأشياء التي تقول فرضيتك إنها موجودة. ولأن~$\OLId{latin-looped}{A}$ تحتوي شيئًا واحدًا على
الأقل، توجد أشياء قد يشير إليها ذلك الاسم. وقد لا تستطيع بالطبع
تعيين أحدها أو وصفه أكثر (سوى أنه~$\in \OLId{latin-looped}{A}$). لكنك تستطيع، لغرض
البرهان، أن تتظاهر بأنك عينته وأن تسميه. ومن المهم أن تختار اسمًا لم
تستخدمه بالفعل (ولم يرد في فرضياتك)، وإلا فقد تختل الأمور. وتشير إلى
ذلك في برهانك بالانتقال من «لبعض~$\OLId{latin-ordinary}{x}$،~$\OLId{latin-ordinary}{x} \in \OLId{latin-looped}{A}$» إلى «ليكن
$\OLId{latin-ordinary}{a} \in \OLId{latin-looped}{A}$». ويمكنك الآن أن تستدل بشأن~$\OLId{latin-ordinary}{a}$، وتستخدم فرضيات أخرى، وهلم
جرًا، إلى أن تبلغ نتيجة~$\OLId{latin-ordinary}{p}$. فإذا لم تعد~$\OLId{latin-ordinary}{p}$ تذكر~$\OLId{latin-ordinary}{a}$، كانت~$\OLId{latin-ordinary}{p}$
مستقلة عن الافتراض~$\OLId{latin-ordinary}{a} \in \OLId{latin-looped}{A}$، وتكون قد بينت أنها تلزم فقط من الافتراض
«لبعض~$\OLId{latin-ordinary}{x}$،~$\OLId{latin-ordinary}{x} \in \OLId{latin-looped}{A}$».

\begin{prop}
إذا كان~$\OLId{latin-looped}{A} \neq \emptyset$، فإن~$\OLId{latin-looped}{A} \cup \OLId{latin-looped}{B} \neq \emptyset$.
\end{prop}

\begin{proof}
  افرض~$\OLId{latin-looped}{A} \neq \emptyset$. فيوجد، من ثم، بعض~$\OLId{latin-ordinary}{x}$ بحيث~$\OLId{latin-ordinary}{x} \in \OLId{latin-looped}{A}$.
  \begin{quote}
    ليس ما فعلناه أولًا إلا إعادة عبارة الفرض. فإن~$\OLId{latin-looped}{A} \neq \emptyset$
    تستبطن دعوى وجودية تظهر بعد بسط جملة من التعريفات. فتعريف~$=$
    يقضي بأن~$\OLId{latin-looped}{A} = \emptyset$ إذا وفقط إذا كان، لكل~$\OLId{latin-ordinary}{x}$: إذا
    $\OLId{latin-ordinary}{x} \in \OLId{latin-looped}{A}$ فإن~$\OLId{latin-ordinary}{x} \in \emptyset$، وإذا~$\OLId{latin-ordinary}{x} \in \emptyset$ فإن
    $\OLId{latin-ordinary}{x} \in \OLId{latin-looped}{A}$. وبنفي الجانبين يصير~$\OLId{latin-looped}{A} \neq \emptyset$ مكافئًا لوجود
    $\OLId{latin-ordinary}{x} \in \OLId{latin-looped}{A}$ بحيث~$\OLId{latin-ordinary}{x} \notin \emptyset$، أو لوجود~$\OLId{latin-ordinary}{x} \in \emptyset$
    بحيث~$\OLId{latin-ordinary}{x} \notin \OLId{latin-looped}{A}$. ولما كان~$\OLId{latin-ordinary}{x} \in \emptyset$ ممتنعًا لأي~$\OLId{latin-ordinary}{x}$،
    بطل الطرف الثاني، واختصر اقتران «$\OLId{latin-ordinary}{x} \in \OLId{latin-looped}{A}$ و
    $\OLId{latin-ordinary}{x} \notin \emptyset$» إلى~$\OLId{latin-ordinary}{x} \in \OLId{latin-looped}{A}$. فثبت أن
    $\OLId{latin-looped}{A} \neq \emptyset$ إذا وفقط إذا كان~$\OLId{latin-ordinary}{x} \in \OLId{latin-looped}{A}$ لبعض~$\OLId{latin-ordinary}{x}$؛ وهذه
    دعوى وجودية. والآن نضع اسمًا لأحد !!{element}s~$\OLId{latin-looped}{A}$:
  \end{quote}
  ليكن~$\OLId{latin-ordinary}{a} \in \OLId{latin-looped}{A}$.
  \begin{quote}
    سمّينا واحدًا من الأشياء التي~$\in \OLId{latin-looped}{A}$، وسنجري الاستدلال على~$\OLId{latin-ordinary}{a}$
    من غير أن نحمّله فرضًا زائدًا على~$\OLId{latin-ordinary}{a} \in \OLId{latin-looped}{A}$:
  \end{quote}
  لما كان~$\OLId{latin-ordinary}{a} \in \OLId{latin-looped}{A}$، كان~$\OLId{latin-ordinary}{a} \in \OLId{latin-looped}{A} \cup \OLId{latin-looped}{B}$ بتعريف~$\cup$. ومن ثم
  يوجد بعض~$\OLId{latin-ordinary}{x}$ بحيث~$\OLId{latin-ordinary}{x} \in \OLId{latin-looped}{A} \cup \OLId{latin-looped}{B}$، أي~$\OLId{latin-looped}{A} \cup \OLId{latin-looped}{B} \neq \emptyset$.
  \begin{quote}
    انتقلنا أخيرًا من «$\OLId{latin-ordinary}{a} \in \OLId{latin-looped}{A} \cup \OLId{latin-looped}{B}$» إلى «لبعض~$\OLId{latin-ordinary}{x}$،
    $\OLId{latin-ordinary}{x} \in \OLId{latin-looped}{A} \cup \OLId{latin-looped}{B}$». ولما خلت العبارة الثانية من~$\OLId{latin-ordinary}{a}$، علمنا أن
    «لبعض~$\OLId{latin-ordinary}{x}$،~$\OLId{latin-ordinary}{x} \in \OLId{latin-looped}{A} \cup \OLId{latin-looped}{B}$» لازمة عن «لبعض~$\OLId{latin-ordinary}{x}$،~$\OLId{latin-ordinary}{x} \in \OLId{latin-looped}{A}$»
    وحدها؛ ومؤدى ذلك~$\OLId{latin-looped}{A} \cup \OLId{latin-looped}{B} \neq \emptyset$.
  \end{quote}
\end{proof}

والأحوط أن تميز المتغيرات المقيدة مثل «$\OLId{latin-ordinary}{x}$» من الأسماء الموضوعة
للشهود مثل~$\OLId{latin-ordinary}{a}$، كما صنعنا. غير أن العمل الجاري كثيرًا ما يستعمل~$\OLId{latin-ordinary}{x}$
للأمرين اختصارًا، على هذا الوجه:
\begin{quote}
افرض~$\OLId{latin-looped}{A} \neq \emptyset$، أي يوجد~$\OLId{latin-ordinary}{x} \in \OLId{latin-looped}{A}$. وبتعريف~$\cup$،
$\OLId{latin-ordinary}{x} \in \OLId{latin-looped}{A} \cup \OLId{latin-looped}{B}$؛ ومن ثم~$\OLId{latin-looped}{A} \cup \OLId{latin-looped}{B} \neq \emptyset$.
\end{quote}
فإذا اختُصر على هذا النحو لزم الاحتراز، وجعل~$\OLId{latin-ordinary}{x}$ و$\OLId{latin-ordinary}{y}$ رمزين مختلفين
لدعويين وجوديتين مختلفتين. ولذلك فإن الكلام الآتي \emph{ليس} برهانًا
صحيحًا على «إذا كان~$\OLId{latin-looped}{A} \neq \emptyset$ و$\OLId{latin-looped}{B} \neq \emptyset$، فإن
$\OLId{latin-looped}{A} \cap \OLId{latin-looped}{B} \neq \emptyset$»؛ بل العبارة نفسها كاذبة.
\begin{quote}
افرض~$\OLId{latin-looped}{A} \neq \emptyset$ و$\OLId{latin-looped}{B} \neq \emptyset$. فيكون لبعض~$\OLId{latin-ordinary}{x}$
$\OLId{latin-ordinary}{x} \in \OLId{latin-looped}{A}$، ويكون لبعض~$\OLId{latin-ordinary}{x}$ كذلك~$\OLId{latin-ordinary}{x} \in \OLId{latin-looped}{B}$. وبما أن~$\OLId{latin-ordinary}{x} \in \OLId{latin-looped}{A}$ و
$\OLId{latin-ordinary}{x} \in \OLId{latin-looped}{B}$، كان~$\OLId{latin-ordinary}{x} \in \OLId{latin-looped}{A} \cap \OLId{latin-looped}{B}$ بتعريف~$\cap$؛ ومن ثم
$\OLId{latin-looped}{A} \cap \OLId{latin-looped}{B} \neq \emptyset$.
\end{quote}
فأين وقع الانتقال الفاسد، ولِمَ لا تثبت منه النتيجة؟

\end{document}
